\section{Background and Threat Model}
\label{sec:background}

\subsection{Agent Skill Packages}

An Agent Skill is treated as a versioned package containing a primary skill
document and any referenced instructions, scripts, configuration, manifests,
and resources.  The package boundary is the unit of analysis.  A single
document may mix user-facing descriptions with operative instructions, so
content roles are assigned at section or paragraph granularity rather than by
filename alone.

\subsection{Threat Model}

The adversary controls every untrusted byte in the package and may distribute
behavior across artifacts.  The adversary knows the analysis strategy and may
use paraphrase, aliases, indirection, string construction, dynamic imports,
reflection, simple encoding, or opaque binaries.  The detector and its policy
and rule versions are trusted.  The detector neither installs dependencies nor
imports, invokes, or executes target-supplied content, and it does not contact
endpoints named by the package.

The method does not claim to recover code obtained only at runtime, behavior
inside unavailable dependencies, or semantics hidden in unrecoverable encrypted
or binary content.  Such uncertainty is represented in the output.  In
particular, \textsc{Pass} is not a proof of safety.

\subsection{Design Goals}

The detector should (G1) remain zero execution; (G2) preserve independence
between high-level function and internal behavior; (G3) connect cross-file
evidence; (G4) reduce false positives on legitimate sensitive operations;
(G5) expose uncertainty instead of silently treating missing analysis as
benign; and (G6) return evidence that a reviewer can locate in the original
package.
